\section{Proof of Corollary~\ref{cor:local}}\label{app:localized}
We use the common-mesh variant of Algorithm~1 (mesh $h_j=s2^{-j}$ in all
coordinates, corrections $L_iw_i(v)^2/8$, two-point grids for coordinates with
$L_i=0$) under point growth with $8\kappa\theta^2\le1$. The proofs of
Lemmas~\ref{lem:inv} and~\ref{lem:states}, written with this mesh, give at stage
$j$ with center $c=y_{j-1}$ and corrected minimizer $y_j$:
\begin{enumerate}[label=(G\arabic*)]
\item $D(z)\le\frac L4nh_j^2+\frac{L\theta^2}2(\norm{z-x^*}^2+\norm{c-x^*}^2)$ for
every grid point $z$;
\item $\norm{y_j-x^*}^2\le\kappa nh_j^2$ and $\norm{c-x^*}^2\le4\kappa nh_j^2$;
\item $U-\OPT\le F(y_j)-\OPT\le D(y_j)\le\frac78Lnh_j^2$;
\item for $i\in I_C$ with $L_i>0$ the filtered interval lies within
$5\sqrt{n\kappa}\,h_j$ of $(y_j)_i$, and within $1+5\sqrt{n\kappa}\,h_j$ for
$i\in I_Z$.
\end{enumerate}
(The sharper constants $\frac9{16}$ and $4.2$ also hold.) For a coordinate with
$L_i=0$ the correction vanishes and (G1)--(G3) hold verbatim.

\begin{proof}
Note that $S\cap C=\emptyset$: for $i\in S$, Lemma~\ref{lem:growthcert}(b)
gives $H_{ii}\ge2g>0$. Similarly, for $i\in C$ the function
$t\mapsto F(x^*+te_i)$ is concave along the coordinate line through $x^*$
and $F(x^*+te_i)-\OPT\ge gt^2>0$ for feasible $t\ne0$, so $x^*_i$ is an
endpoint of $X_i$; also $B_i=B^\circ_i=X_i$, because a coordinate with a
single grid interval is never narrowed. Fix a stage $j$ with $h_j\le h^*$.

\emph{Node exclusion.} Let $i\in Z\cup C$ and let $v\ne x^*_i$ be a node of
$G_i$; put $\delta=|v-x^*_i|$, so $\delta\ge1$ for $i\in Z$ and
$\delta=u_i-l_i\ge\delta_C$ for $i\in C$. Let $z$ be a grid point with
$z_i=v$ and $Q(z)=m_i(v)$. By growth and (G1), (G2), and
$L\theta^2/2\le g/16$,
\[
m_i(v)-\OPT\ge g\norm{z-x^*}^2-D(z)\ge\tfrac{15}{16}g\delta^2
-\Bigl(\frac L4+\frac{g\kappa}4\Bigr)nh_j^2 .
\]
With (G3) and $L\le\kappa g$ this gives
$m_i(v)-U\ge g\bigl(\frac{15}{16}\delta^2-\frac{11}8\kappa nh_j^2\bigr)>0$
whenever $h_j^2<\frac{15}{22}\delta^2/(\kappa n)$, which holds since
$h_j\le h^*$ and $\frac12<\frac45<\sqrt{15/22}$. Hence for $i\in C$
the rule sets $B'_i=\{x^*_i\}$ (the node $x^*_i$ is never excluded, because
$m_i(x^*_i)\le\OPT\le U$ by Lemma~\ref{lem:node}). For $i\in Z$, all
nodes other than $x^*_i$ have $m_i(v)>U$. Moreover, by (G2),
$|c_i-x^*_i|\le2\sqrt{\kappa n}\,h_j\le1$, and every
integer step taken within distance $5$ of $c_i$ equals
$\max\{1,\lfloor h_j+\theta t\rfloor\}=1$ because
$h_j+\theta t\le\frac12+\frac54<2$. So the grid contains every integer of
$B^\circ_i$ within distance $4$ of $x^*_i$, all grid intervals meeting
$[x^*_i-4,x^*_i+4]$ have length one, and by (G4) and (G2) the hull $B_i$
lies within $1+6\sqrt{n\kappa}\,h_j\le4$ of $x^*_i$. Thus $K_i\cap B_i=\{x^*_i\}$
and $B'_i=\{x^*_i\}$.

\emph{The remaining coordinates.} Hence $W\subseteq S\cup A$, and
$x^*\in B'$ because $x^*\in B$ by (F). For $i\in S\cap W$ we have
$\partial_iF(x^*)=0$, so case (i), (ii) or (iii) applies with $\mu_i=0$. For
$i\in A\cap W$, $x^*_i$ is an endpoint of $X_i$, hence of $B'_i=B_i$, and
the first-order conditions give the sign required by case (i) or (ii); by
(G4) the length of $B'_i$ is at most $10\sqrt{n\kappa}\,h_j$, so
$\mu_i\ge\lambda_A/(10\sqrt{n\kappa}\,h_j)\ge\Gamma/2$.

\emph{Positive semidefiniteness.} Let $A'=A\cap W$ and
$N=H_{WW}+2\operatorname{diag}(\mu_W)$. If $A'=\emptyset$, then
$N=H_{SS}\succeq2gI$ by Lemma~\ref{lem:growthcert}(b) (all of $S$ lies in
$W$, since a continuous coordinate with $L_i>0$ keeps a hull of positive
length). Otherwise $H_{SS}\succ0$, and $N\succeq0$ if the Schur complement
$H_{A'A'}+2\operatorname{diag}(\mu_{A'})-H_{A'S}H_{SS}^{-1}H_{SA'}$ is
positive semidefinite. Its smallest eigenvalue is at least
$2\min_{A'}\mu_i-\norm{H_{AA}}_2-\norm{H_{SA}}_2^2/(2g)\ge0$, using
$\norm{H_{SS}^{-1}}_2\le1/(2g)$ and that submatrices have smaller norms.
Finally $F(x^*)=\OPT\le U$, so all hypotheses of
Proposition~\ref{prop:local} hold. The stage count follows from
$h_j=s2^{-j}$.
\end{proof}
